\documentclass[10pt,twocolumn,letterpaper]{article}

\usepackage{cvpr}
\usepackage{times}
\usepackage{epsfig}
\usepackage{graphicx}
\usepackage{amsmath}
\usepackage{amssymb}
\usepackage{booktabs}
\usepackage[breaklinks=true,bookmarks=false]{hyperref}

\cvprfinalcopy
\def\cvprPaperID{****}
\def\httilde{\mbox{\tt\raisebox{-.5ex}{\symbol{126}}}}
\setcounter{page}{1}

\begin{document}

\title{NavLM Without a Compass: Self-Supervised Visual Walking Directions from Zurich Walking Tours}

\author{Author One\\
Stanford University\\
{\tt\small author1@stanford.edu}
\and
Author Two\\
Stanford University\\
{\tt\small author2@stanford.edu}
\and
Author Three\\
Stanford University\\
{\tt\small author3@stanford.edu}
}

\maketitle

\begin{abstract}
We study a practical version of visual navigation assistance: given a phone photo of the current scene, a target destination, and a route summary computed from a street graph, can a vision-language model produce the correct next walking action without access to a compass? We build a self-supervised pipeline over Zurich walking-tour videos that aligns video frames with GPS, Street View, OpenStreetMap (OSM), and a curated set of 21 tourist destinations, then converts those alignments into instruction-tuning data for a four-way action decision: \textit{continue ahead}, \textit{turn left}, \textit{turn right}, or \textit{turn around}.

Our main contribution is an end-to-end data engine and experimental design for testing whether heading can be removed from the student model's input while preserving useful navigation behavior. The repo implements three prompt conditions: \textit{given} (heading provided), \textit{derived} (heading hidden but inferred in chain-of-thought), and \textit{implicit} (heading hidden and not stated numerically). The current pipeline yields 1,219 matched frames, 3,657 routed frame-destination pairs, and per-variant instruction-tuning datasets with 2,472 / 1,656 / 2,344 training examples for given / derived / implicit. Preliminary zero-shot smoke tests on 16 held-out examples show a large gap between heading-given and compass-free settings, motivating the full fine-tuning sweep that the codebase is prepared to run.
\end{abstract}

\section{Introduction}
Navigation assistants are now common, but most systems still depend on map UIs, turn-by-turn text, or explicit device orientation. For a pedestrian standing in an unfamiliar city, a more natural interface is: take a photo, name the destination, and receive a short spoken next-step instruction grounded in what is visible. This problem combines visual place understanding, local orientation, route-following, and language generation.

This project focuses on one concrete question: \textit{can a vision-language model issue the correct next walking action from a single street photo without being given a compass heading?} The challenge is not route planning alone; the route can be computed from OSM. The challenge is to align the route direction with the walker's current visual orientation. Our hypothesis is that in many urban scenes, a model can recover enough orientation from appearance cues such as church spires, river-bank geometry, street axes, tram tracks, storefronts, and landmark visibility to choose the correct first action even when numeric heading is withheld.

The input to our algorithm is a tuple
\begin{equation}
x = (I, d, r, L, h_{\text{opt}}),
\end{equation}
where $I$ is a phone image, $d$ is one of 21 curated Zurich attractions, $r$ is a route summary derived from the OSM walking graph, $L$ is a list of visible landmarks, and $h_{\text{opt}}$ is an optional heading value. The output is a short natural-language instruction whose first action verb must be one of four classes: \texttt{continue ahead}, \texttt{turn left}, \texttt{turn right}, or \texttt{turn around}.
The \textit{given} variant includes $h_{\text{opt}}$, while the \textit{derived} and \textit{implicit} variants hide it.

Our system is built in two stages. First, we construct self-supervised training data by aligning Zurich walking-tour frames with Street View and GPS, then matching GPS-side and VLM-side landmark evidence. Second, we fine-tune Qwen2.5-VL-7B with LoRA adapters on teacher-labeled navigation examples generated by Gemini Pro 2.5. The resulting report is therefore both a machine learning study and a systems paper: the main technical work lies in building reliable labels from weak multimodal supervision rather than collecting a hand-annotated dataset from scratch.

\section{Related Work}
Our project sits at the intersection of vision-language navigation, visual geo-localization, and parameter-efficient adaptation of multimodal models.

\textbf{Vision-language navigation.} Classical VLN benchmarks such as Room-to-Room (R2R) and Room-across-Room (RxR) study instruction following in indoor or simulator-based environments \cite{anderson2018vision,ku2020room}. These datasets established the importance of grounding language in visual observations, but they typically assume embodied agents in discrete environments with provided trajectories. Touchdown is especially relevant to our setting because it studies outdoor street-level navigation using real imagery \cite{chen2019touchdown}. Our task differs in that the model is not asked to execute a full trajectory from scratch. Instead, it must pick the correct \textit{next action} from a single current image plus a graph-derived route summary.

\textbf{Visual geo-localization and place recognition.} Outdoor navigation depends on robust place recognition and orientation cues. NetVLAD and later large-scale localization systems show how learned image descriptors can align ground images with map or reference imagery \cite{arandjelovic2016netvlad,sarlin2019coarse}. DINOv2 offers strong general-purpose visual features and is used in our pipeline to align video frames to Street View crops \cite{oquab2023dinov2}. Unlike pure geo-localization, however, our goal is not only to identify a place but also to determine the action that aligns the walker's facing direction with a route.

\textbf{Multimodal instruction and city-scale grounding.} Recent large multimodal models can reason over images and structured prompts, making them attractive as teachers or students for grounded navigation. Gemini is used here as a teacher model for data synthesis \cite{gemini2024}. Qwen2.5-VL is the student architecture we fine-tune \cite{qwen25vl2025}. Our formulation is narrower than open-ended embodied-agent systems: we deliberately constrain output to four verbs so we can measure directional correctness precisely.

\textbf{Parameter-efficient adaptation.} We train only small low-rank adapters rather than full model weights. LoRA and QLoRA show that low-rank updates on quantized base models can retain strong downstream performance at much lower memory cost \cite{hu2022lora,dettmers2023qlora}. This is important for our project because the same base VLM is trained under three prompt conditions and multiple adapter ranks.

\textbf{Maps, open infrastructure, and tooling.} Our labels rely on OSM and its walking graph, which remain a practical foundation for city-scale academic projects \cite{haklay2008openstreetmap}. We implement the model stack in PyTorch and Hugging Face Transformers \cite{paszke2019pytorch,wolf2020transformers}.

Relative to prior work, our contribution is not a new backbone architecture. Instead, we provide a self-supervised recipe for converting unlabelled walking-tour videos into a controlled urban next-action benchmark, and we use it to isolate the role of heading information in multimodal pedestrian guidance.

\section{Methods}
\subsection{Overview}
The repo's ``Attempt 2'' pipeline was motivated by a failure in an earlier version of the project: using nearest OSM place names as \textit{destinations} caused the model to learn street names instead of tourist-relevant endpoints. The revised pipeline separates \textit{source location} from \textit{destination}. A frame's GPS may lie on a street, but the destination is sampled from a curated 21-attraction vocabulary.

The full pipeline has five stages:
\begin{enumerate}
\item Build GPS-side and VLM-side candidate landmark lists for each frame.
\item Match frames whose GPS-derived and VLM-derived names coincide.
\item Estimate heading, snap GPS to the walking graph, and define routable destination targets.
\item Generate deterministic route-derived ground-truth verbs for sampled frame-destination pairs.
\item Distill Gemini-generated responses into variant-specific SFT datasets and fine-tune Qwen2.5-VL-7B adapters.
\end{enumerate}

\subsection{Frame matching and destination construction}
Let $f_i$ denote a video frame. For each accepted frame, the GPS-side module builds three sets: nearby curated attractions, nearby OSM landmarks, and nearby OSM POIs. In parallel, a VLM-side module extracts attraction, landmark, and POI names from the model's visible landmark list and place guess. A frame is retained if at least one name coincides across the GPS-side and VLM-side unions through exact match, substring match, or informative word overlap.

This step is deliberately permissive. We do not require exact string identity because tourist landmarks appear under many aliases (e.g., ``St. Peter Church'' vs.\ ``St. Peter''). The result is a matched cohort that is semantically aligned even when raw naming conventions differ.

\subsection{Heading and route-derived action labels}
The current heading is estimated from DINOv2 similarities over compass crops. If the similarity gap between the top-1 and top-2 orientations is large, the system trusts the top-1 angle; otherwise it uses a cosine-weighted blend of the top two. For a routed frame-destination pair, the OSM graph provides the first route bearing $b$ from the current snapped node to the next node on the shortest path.

The ground-truth action is chosen by applying each candidate verb's angle offset and minimizing angular error to the route bearing:
\begin{equation}
a^\star = \arg\min_{a \in \mathcal{A}}
\left|
\angle \left((h + \Delta(a)) \bmod 360,\; b\right)
\right|,
\end{equation}
where $\Delta(\texttt{continue ahead}) = 0^\circ$, $\Delta(\texttt{turn left}) = -90^\circ$, $\Delta(\texttt{turn right}) = 90^\circ$, and $\Delta(\texttt{turn around}) = 180^\circ$. This gives deterministic labels without relying on teacher free-form preferences.

\subsection{Three prompt conditions}
We compare three student prompting conditions:
\begin{itemize}
\item \textbf{Given:} the student receives the numeric heading and route bearing.
\item \textbf{Derived:} the heading is hidden from the student; the model is asked to infer it from the image and then choose a verb.
\item \textbf{Implicit:} the heading is hidden and the model is asked to reason visually without emitting a numeric heading.
\end{itemize}

Gemini Pro 2.5 labels each routed example three separate times, once per condition. The teacher always has access to the true heading so it can compute the correct action, but its chain-of-thought style is conditioned to match the student condition. This creates supervision with asymmetric information: the student may not receive the heading even though the teacher used it internally.

\subsection{Model and optimization}
The student is Qwen2.5-VL-7B-Instruct, trained with 4-bit NF4 quantization and LoRA adapters on the attention projection layers. We sweep LoRA rank $r \in \{4,8,16\}$ while keeping the learning rate fixed at $2\times10^{-4}$ and using $2$ epochs, batch size $1$, and gradient accumulation $8$.

An important implementation detail is that the training loss is masked to the assistant turn only. The system and user prompt tokens are excluded from cross-entropy, so optimization focuses on generating the teacher's \texttt{<thinking>} and \texttt{<answer>} tokens rather than memorizing repeated prompt boilerplate. This matters because all three variants share highly regular prompt structure.

\subsection{What we implemented vs.\ what we reused}
This project builds on an earlier internal NavLM codebase and prior attempt artifacts, but the key Attempt 2 components are newly written in the current repo. Reused components include the Street View crawl, DINOv2 matching cache, OSM walking graph, and an earlier road-snapping module. New code in \texttt{src/a2\_*.py} implements the curated destination vocabulary, GPS/VLM agreement filter, improved heading estimator, destination-target construction, route labeling, teacher prompting, SFT conversion, Modal training/evaluation harnesses, and scoring pipeline.

\section{Dataset and Features}
\subsection{Source data}
The raw source data consists of eight unlabelled 4K Zurich walking-tour videos listed in the repo, together with downloaded Google Street View reference crops and OSM map data. The processing region covers Zurich's old town and nearby tourist areas. The project also maintains a held-out video identifier in configuration, although the current Attempt 2 SFT splits are random per variant rather than strictly video-held-out.

\subsection{Constructed dataset sizes}
The current pipeline produces the following intermediate and final datasets:
\begin{itemize}
\item 15,053 DINO-accepted video frames with GPS-side metadata.
\item 4,891 VLM-scanned frames with VLM-side landmark metadata.
\item 1,219 matched frames where GPS-side and VLM-side evidence coincide.
\item 89 matched panoramas and 21 curated destination attractions.
\item 3,657 routed frame-destination pairs (three sampled destinations per matched frame).
\end{itemize}

Teacher annotation is nearly complete but not perfectly symmetric across variants in the current snapshot. The repo presently contains 3,657 given, 2,866 derived, and 3,605 implicit annotated examples. After filtering for well-formed teacher outputs, the SFT splits are:
\begin{itemize}
\item \textbf{Given:} 2,472 train / 309 val / 309 test.
\item \textbf{Derived:} 1,656 train / 206 val / 206 test.
\item \textbf{Implicit:} 2,344 train / 293 val / 293 test.
\end{itemize}

\subsection{Features and supervision signals}
Our model uses both visual and structured inputs. Visual features come from the RGB image itself; alignment features for data construction come from DINOv2 similarity to Street View crops. Structured features include destination identity, route distance, route first-segment bearing, visible landmark names, and optionally the heading. The crucial supervision signal is the deterministic first action extracted from the street-graph shortest path and the estimated heading.

This makes the benchmark different from generic captioning or geolocation. The image alone is not sufficient, and the structured route summary alone is not sufficient. The model must combine them.

\begin{figure}[t]
\centering
\fbox{\parbox{0.92\linewidth}{\vspace{1.2in}\centering
Replace with screenshots from \texttt{viz/a2\_vlmagreed.html}\\
and \texttt{viz/a2\_mapped\_GPS\_spot.html}.\\
Show matched-frame QC and the spatial distribution of the 89 matched panoramas.
\vspace{1.2in}}}
\caption{Suggested qualitative dataset figure. The repo already contains HTML visualizations for matched-frame inspection and map coverage; these should be exported to static figures for the final paper.}
\label{fig:data}
\end{figure}

\section{Experiments, Results, and Discussion}
\subsection{Experimental setup}
We evaluate 12 conditions in total: three zero-shot conditions and nine trained conditions (three prompt variants times three LoRA ranks). The current repo snapshot contains complete code for all 12 conditions, but only preliminary zero-shot smoke-test outputs for the new Attempt 2 setup are available locally. Full trained results are therefore future work in the narrow sense of this draft, but the paper can still present the benchmark, methodology, and preliminary trends.

We report four metrics:
\begin{equation}
\text{format\_pass} = \mathbb{1}[\texttt{<thinking>} \land \texttt{<answer>} \land \hat{a}\neq \varnothing],
\end{equation}
\begin{equation}
\text{direction\_pass} = \mathbb{1}[\hat{a} = a^\star],
\end{equation}
\begin{equation}
\text{PASS} = \text{format\_pass} \cdot \text{direction\_pass},
\end{equation}
and for derived conditions only, a heading inference accuracy that checks whether the model's claimed heading is within $22.5^\circ$ of the true heading when such a value is emitted.

\subsection{Preliminary quantitative results}
Table~\ref{tab:zs} summarizes the current zero-shot smoke-test results on 16 examples per condition. These numbers should be interpreted as a sanity check rather than a final benchmark, but they already illustrate the project's central difficulty.

\begin{table}[t]
\centering
\small
\begin{tabular}{lcccc}
\toprule
Condition & $n$ & Format & Direction & PASS \\
\midrule
ZS given & 16 & 100.0 & 62.5 & 62.5 \\
ZS derived & 16 & 100.0 & 18.8 & 18.8 \\
ZS implicit & 16 & 100.0 & 12.5 & 12.5 \\
\bottomrule
\end{tabular}
\caption{Preliminary zero-shot smoke-test results from the current repo snapshot. All numbers are percentages.}
\label{tab:zs}
\end{table}

The gap between \textit{given} and the compass-free variants is large. This is precisely the behavior we expected: once numeric heading is hidden, the model must infer orientation from visual cues and city knowledge, which is much harder than subtracting angles. The derived variant performs slightly better than implicit in this tiny sample, suggesting that explicitly asking the model to recover heading may help, but the sample is too small to support strong claims.
Appendix~A includes a ready-to-fill full-results table for the eventual 12-condition sweep, along with additional qualitative templates that can be populated once more outputs are available.

\subsection{Teacher quality and dataset quality}
Teacher quality is also uneven across variants. In the current annotation snapshot, the number of examples with overall teacher PASS is 3,201 / 2,061 / 2,632 for given / derived / implicit. The lower count for derived reflects the longer and more failure-prone reasoning format. This creates a trade-off between richer supervision and annotation reliability. In our current training setup we keep all rows with valid format by default, not only rows with correct teacher direction, in order to preserve dataset size.

\subsection{Qualitative analysis and expected figures}
The repo already contains useful qualitative artifacts:
\begin{itemize}
\item \texttt{viz/a2\_vlmagreed.html}: agreement between GPS-side and VLM-side matching.
\item \texttt{viz/a2\_route\_gt.html}: route and ground-truth action visualization.
\item \texttt{viz/a2\_viz\_sft.html}: spot-checks of generated SFT examples.
\item \texttt{viz/a2\_thin\_attractions.html}: failure analysis for thin-coverage destinations.
\end{itemize}

For the final submission, we would convert these into two or three static figures: one showing matched-frame QC, one showing route/action examples, and one showing typical success/failure cases. The most important failure modes to highlight are: (1) ambiguous orientation on open squares or riverfront scenes, (2) long thin destinations such as riverbanks and streets, and (3) sparse destinations with very few matched frames.

\begin{figure}[t]
\centering
\fbox{\parbox{0.92\linewidth}{\vspace{1.0in}\centering
Replace with screenshots from \texttt{viz/a2\_route\_gt.html}\\
and \texttt{viz/a2\_viz\_sft.html}.\\
Show one success case and one compass-free failure case.
\vspace{1.0in}}}
\caption{Suggested qualitative result figure. Use one routed example where the model succeeds with strong visual cues and one where orientation is ambiguous.}
\label{fig:qual}
\end{figure}

\subsection{Discussion}
Even before the full fine-tuning sweep, the project yields several lessons. First, destination design matters: replacing noisy nearest-street names with a curated attraction vocabulary is essential. Second, agreement between independent supervision sources (GPS/OSM and VLM landmark recognition) is a practical way to build a high-precision cohort from noisy video data. Third, the problem is not simply ``geolocation''; it is \textit{orientation-conditioned action prediction}. The strong zero-shot drop from given to derived/implicit suggests that heading is the core bottleneck and that the final scientific question is well posed.

We do not yet claim that the compass-free thesis holds. The correct claim at this stage is narrower: the repo now supports a controlled experiment that can test it.

\section{Conclusion and Future Work}
We presented a self-supervised pipeline for learning image-grounded pedestrian directions in Zurich from walking-tour videos, Street View, and OSM. The key idea is to convert weak multimodal alignments into a next-action benchmark where a multimodal model must reconcile what it sees with a graph-derived route. The current codebase and constructed datasets are substantial: 1,219 matched frames, 3,657 routed examples, variant-specific SFT splits, and a full LoRA training/evaluation harness.

Our preliminary results show that zero-shot performance is much stronger when heading is supplied than when it must be inferred. This validates the importance of the problem formulation. The next step is to complete the 12-condition sweep and determine whether fine-tuning on teacher-generated compass-free reasoning narrows the gap enough to support our main hypothesis.

If we had more time and compute, we would expand beyond one city, replace single-step evaluation with multi-step rollouts or progress-based navigation metrics, and perform a human study on instruction usefulness. We would also investigate richer destination formulations and stronger orientation supervision from map geometry or multi-frame context.

\appendix
\section{Appendix: Visualization Plan and Fill-In Tables}
This appendix is intentionally drafted as a working scaffold. The goal is to make it easy to drop in screenshots, examples, and final quantitative values after the main experiments finish.

\subsection{DINOv2 Matching Success and Failure Cases}
One useful appendix figure is a side-by-side panel built from \texttt{viz/a2\_vlmagreed.html}. This figure can show: (1) a \textbf{success case}, where the query frame, top Street View crop, and GPS/VLM landmark lists are all semantically consistent; and (2) a \textbf{failure or borderline case}, where DINOv2 produces a visually similar but geographically ambiguous match, or where heading estimation becomes unstable because multiple compass crops have similar cosine scores.

When discussing these examples, it is helpful to point out three specific failure patterns:
\begin{itemize}
\item \textbf{Lookalike facades and old-town streets.} Zurich's old town contains many narrow cobblestone streets and repetitive building fronts, so visually similar locations can be confused even when local texture looks plausible.
\item \textbf{Weak heading margins.} In some cases the best and second-best compass crops have similar cosine similarity, which makes the heading estimate less stable and increases downstream verb errors.
\item \textbf{Good place match, weak orientation match.} A frame can be roughly geo-localized correctly while still failing to recover the camera's facing direction precisely enough for the correct first action.
\end{itemize}

\begin{figure}[t]
\centering
\fbox{\parbox{0.92\linewidth}{\vspace{1.0in}\centering
Appendix figure placeholder:\\
left = DINOv2 success case from \texttt{viz/a2\_vlmagreed.html}\\
right = DINOv2 failure / ambiguous-heading case from \texttt{viz/a2\_vlmagreed.html}
\vspace{1.0in}}}
\caption{Appendix placeholder for DINOv2 matching analysis. Suggested discussion: query frame, top Street View match, cosine margin, GPS/VLM agreement, and why the example succeeds or fails.}
\label{fig:appendix_dino}
\end{figure}

\subsection{OSM Routing Examples}
Another useful appendix figure is a route visualization from \texttt{viz/a2\_route\_gt.html}. A strong appendix treatment would include one \textbf{point-target} example and one \textbf{multi-target} example. The point-target example can emphasize a short, unambiguous route to a church or square, while the multi-target example can emphasize why network distance matters for long features such as the Limmat river or Bahnhofstrasse.

The appendix discussion can explain:
\begin{itemize}
\item how the walker's snapped node is selected;
\item how the OSM shortest path determines the first-segment route bearing;
\item how the route bearing and estimated heading combine into the deterministic ground-truth verb;
\item why multi-target routing is better than using a single centroid for streets or riverfront destinations.
\end{itemize}

\begin{figure}[t]
\centering
\fbox{\parbox{0.92\linewidth}{\vspace{1.0in}\centering
Appendix figure placeholder:\\
OSM routing examples from \texttt{viz/a2\_route\_gt.html}\\
showing one point target and one multi target
\vspace{1.0in}}}
\caption{Appendix placeholder for OSM routing examples. Suggested overlays: walker location, heading arrow, shortest path polyline, destination marker, and verb-error table.}
\label{fig:appendix_route}
\end{figure}

\subsection{Teacher and Student Prompt Examples}
The report can also benefit from a prompt appendix showing the asymmetry between teacher and student inputs. A compact presentation is to include one example each for \textit{given}, \textit{derived}, and \textit{implicit}, with the teacher prompt on the left and the student prompt on the right. The most important thing to show is that the teacher always has access to heading, while the student prompt hides it for derived and implicit conditions.

The teacher/student prompt templates are implemented in \texttt{src/a2\_annotate.py}. For the final report, you can either paste shortened prompt excerpts or typeset them as small monospaced blocks.

\begin{figure}[t]
\centering
\fbox{\parbox{0.92\linewidth}{\vspace{0.9in}\centering
Appendix figure placeholder:\\
teacher vs. student prompts for \textit{given}, \textit{derived}, and \textit{implicit}
\vspace{0.9in}}}
\caption{Appendix placeholder for prompt design. Recommended content: one route example, three variants, and explicit highlighting of what information is hidden from the student.}
\label{fig:appendix_prompts}
\end{figure}

\subsection{Full Experiment Table Template}
Table~\ref{tab:full_results_template} is a fill-in template for the full 12-condition sweep. You can keep the zero-shot smoke values already observed, or replace all cells later with the final scored results from \texttt{summary\_table.txt}.

\begin{table*}[t]
\centering
\small
\begin{tabular}{lccccc}
\toprule
Condition & $n$ & Format (\%) & Direction (\%) & PASS (\%) & Heading inference (\%) \\
\midrule
zs-heading-given & -- & -- & -- & -- & n/a \\
zs-heading-derived & -- & -- & -- & -- & -- \\
zs-heading-implicit & -- & -- & -- & -- & n/a \\
trained-heading-given-r4 & -- & -- & -- & -- & n/a \\
trained-heading-given-r8 & -- & -- & -- & -- & n/a \\
trained-heading-given-r16 & -- & -- & -- & -- & n/a \\
trained-heading-derived-r4 & -- & -- & -- & -- & -- \\
trained-heading-derived-r8 & -- & -- & -- & -- & -- \\
trained-heading-derived-r16 & -- & -- & -- & -- & -- \\
trained-heading-implicit-r4 & -- & -- & -- & -- & n/a \\
trained-heading-implicit-r8 & -- & -- & -- & -- & n/a \\
trained-heading-implicit-r16 & -- & -- & -- & -- & n/a \\
\bottomrule
\end{tabular}
\caption{Fill-in template for the full 12-condition experiment matrix. Replace ``--'' with final values after evaluation.}
\label{tab:full_results_template}
\end{table*}

\subsection{Teacher Annotation Quality Table Template}
Table~\ref{tab:teacher_quality_template} can be used to summarize annotation throughput and quality per variant. This is helpful if you want to discuss why the derived variant has a smaller usable training pool.

\begin{table}[t]
\centering
\small
\begin{tabular}{lccc}
\toprule
Variant & Rows & Format-pass & PASS \\
\midrule
given & -- & -- & -- \\
derived & -- & -- & -- \\
implicit & -- & -- & -- \\
\bottomrule
\end{tabular}
\caption{Fill-in template for teacher annotation quality by variant.}
\label{tab:teacher_quality_template}
\end{table}

\subsection{Qualitative Case Summary Table Template}
Table~\ref{tab:qual_case_template} is a compact way to organize the appendix examples once you choose the final screenshots.

\begin{table*}[t]
\centering
\small
\begin{tabular}{p{2.4cm}p{2.6cm}p{2.5cm}p{3.3cm}p{3.8cm}}
\toprule
Case type & File / frame & Destination & What is visible & Why it succeeds or fails \\
\midrule
DINOv2 success & -- & -- & -- & -- \\
DINOv2 failure & -- & -- & -- & -- \\
Routing success & -- & -- & -- & -- \\
Compass-free failure & -- & -- & -- & -- \\
Thin-attraction case & -- & -- & -- & -- \\
\bottomrule
\end{tabular}
\caption{Fill-in template for appendix qualitative examples.}
\label{tab:qual_case_template}
\end{table*}

\section*{Contributions and Acknowledgements}
\textbf{Author contributions.} Replace this paragraph with your actual team information. A recommended format is: \textit{Author One} led dataset construction, GPS/VLM matching, and route generation; \textit{Author Two} implemented model training, evaluation, and ablations; \textit{Author Three} led related work, analysis, figures, and report writing. All authors contributed to problem formulation, experiment design, debugging, and final editing.

\textbf{Acknowledgements.} This project reused components from an earlier internal NavLM codebase, including the Street View crawl, DINOv2 matching cache, and OSM road-snapping infrastructure. We also used OpenStreetMap data, PyTorch, Hugging Face Transformers, and Modal for compute and experiment execution.

\textbf{Generative AI usage statement.} Generative AI tools were used for brainstorming experiment framing, debugging implementation details, improving writing clarity, and drafting report text. All technical claims, code, figures, and final wording were reviewed and verified by the authors, who remain fully responsible for the correctness and originality of the submission.

{\small
\bibliographystyle{ieee}
\bibliography{navlm_refs}
}

\end{document}
